% ================================================================
%  Chapter 7 — Software Testing (KANDA)
%  Knowledge-driven Autonomous Navigation and Decision-making Agent
% ================================================================
\chapter{TESTING}
\label{ch:testing}

This chapter explains how I tested KANDA to make sure it actually works and doesn't crash. Testing was in three layers: unit tests (pure Python logic), integration tests (Pi talking to ESP32), and system tests (the whole robot doing real tasks). This part of the project taught me that hardware testing takes way longer than software testing because you have to physically watch things happen. I also learned that a few simple unit tests catch most of the obvious bugs early.


%-------------------------------------------------------------------
\section{Test Strategy and Environment}
\label{sec:test_strategy}

What I'm testing for. I need to confirm that: (1) the Pi and ESP32 can actually talk to each other without data corruption; (2) the validator never lets through bad commands (wrong action names or speeds outside 0--255); (3) the motors respond correctly when given valid commands; (4) if something fails (network down, API timeout, garbage input), the robot doesn't just freeze, it stops safely; (5) the firmware correctly switches to its own obstacle avoidance when the Pi goes silent. Essentially: can I send good commands and handle bad ones?

Where I tested. Unit tests run on any Linux/macOS/Windows machine with Python 3.10+ and the code checked out. Integration and system tests need the actual hardware: the ESP32, the Raspberry Pi, the sensors, motors, and OLED all wired up correctly. Tests are labeled UT-$n$ (unit—no hardware needed), IT-$n$ (integration—both processors, UART, but not necessarily full autonomy), and ST-$n$ (system—full robot, real cloud API calls, real motion).


%-------------------------------------------------------------------
\section{Unit Testing}
\label{sec:unit_testing}

Unit tests let me check the core logic without wasting time on motors, without needing the network, and without risk of the robot driving into things. I test the validator (does it reject bad commands?) and the telemetry parser (does it understand sensor data?). These tests are quick; they run in seconds and catch obvious bugs. \textbf{Tables~\ref{tab:ut_validator}} and~\textbf{\ref{tab:ut_telemetry}} document the unit test cases.

\subsection[Validator module]{Test suite: safety validator (\textit{safety\_validator.py})}
\label{subsec:ut_validator}

The safety validator is the last line of defense before motors turn on. It needs to reject garbage and clamp speeds. I test it with a built-in list of cases: valid forward, speeds too high (clamped to 255), speeds too low (clamped to 0), made-up action names (mapped to stop), and malformed input (also mapped to stop). Running \textit{python3 safety\_validator.py} prints the results so I can visually verify they're right.


\begin{table}[H]
\centering
\caption[Representative unit test: command validation]{\textbf{Representative unit test: command validation}}
\label{tab:ut_validator}
\begin{tabularx}{\textwidth}{|>{\raggedright\arraybackslash}p{3.1cm}|X|}
\hline
\textbf{Field} & \textbf{Description} \\
\hline
Test ID & UT-01 \\ \hline
Objective & Invalid or malformed commands map to a safe stop; in-range speeds pass; out-of-range speeds clamp. \\ \hline
Input & Dictionaries as in the built-in case list (e.g.\ invalid action \textit{fly} with speed~100; \textit{forward} with speed~300). \\ \hline
Procedure & Execute \textit{validate(cmd)} for each case; compare returned \textit{action} and \textit{speed} to expectations. \\ \hline
Expected & Unknown action or non-dict $\rightarrow$ \textit{stop}, speed 0; forward with 300 $\rightarrow$ forward, 255; forward with $-10$ $\rightarrow$ forward, 0. \\ \hline
Observed & Matches implementation in repository; all listed cases behave as specified. Pass (when script run after code changes). \\
\hline
\end{tabularx}
\end{table}

\subsection[Telemetry parsing]{Test suite: telemetry line parsing (\textit{uart\_bridge.py})}
\label{subsec:ut_telemetry}

Function \textit{\_parse\_telemetry} expects \textit{F:<float> L:<float> R:<float> -> <label>}. Unit verification uses short Python snippets or the module’s standalone send/read smoke test with a loopback or live board.

\begin{table}[H]
\centering
\caption[Representative unit test: telemetry parsing]{\textbf{Representative unit test: telemetry parsing}}
\label{tab:ut_telemetry}
\begin{tabularx}{\textwidth}{|>{\raggedright\arraybackslash}p{3.1cm}|X|}
\hline
\textbf{Field} & \textbf{Description} \\
\hline
Test ID & UT-02 \\ \hline
Objective & Valid lines yield numeric \textit{front}, \textit{left}, \textit{right} and string \textit{action}; malformed lines yield \textit{None}. \\ \hline
Input & \textit{F:45.2 L:30.1 R:80.5 -> FORWARD}; empty string; garbled text. \\ \hline
Procedure & Call parser wrapper or \textit{read\_telemetry} with injected lines. \\ \hline
Expected & First line parses to three positive distances; bad inputs $\rightarrow$ \textit{None}. \\ \hline
Observed & Consistent with \textit{\_parse\_telemetry} implementation. Pass under manual checks. \\
\hline
\end{tabularx}
\end{table}

%-------------------------------------------------------------------
\section{Integration Testing}
\label{sec:integration_testing}

Integration tests require the serial link and firmware running, but may keep the robot on blocks to avoid uncontrolled motion. \textbf{Tables~\ref{tab:it_json}} and~\textbf{\ref{tab:it_telem_up}} specify the integration test cases.

\subsection[UART command round-trip]{UART JSON command round-trip (Pi $\rightarrow$ ESP32)}
\label{subsec:it_uart_cmd}

\begin{table}[H]
\centering
\caption[Integration test: inbound JSON handling]{\textbf{Integration test: inbound JSON handling}}
\label{tab:it_json}
\begin{tabularx}{\textwidth}{|>{\raggedright\arraybackslash}p{3.1cm}|X|}
\hline
\textbf{Field} & \textbf{Description} \\
\hline
Test ID & IT-01 \\ \hline
Objective & ESP32 parses a newline-terminated JSON line and updates mode/OLED without crashing. \\ \hline
Input & Stop then forward JSON lines (one object per line) from the Pi or a serial terminal. \\ \hline
Procedure & Connect UART; send one command at a time; observe serial debug and OLED \textit{Action} field. \\ \hline
Expected & Acknowledged motion or stop; no reset loop; speed within clamped range on scope or duty observation if used. \\ \hline
Observed & Record pass/fail per build; document port name and firmware version. \\
\hline
\end{tabularx}
\end{table}

\subsection[Telemetry uplink]{Telemetry uplink (ESP32 $\rightarrow$ Pi)}
\label{subsec:it_telemetry_up}

\begin{table}[H]
\centering
\caption[Integration test: telemetry reception on the Pi]{\textbf{Integration test: telemetry reception on the Pi}}
\label{tab:it_telem_up}
\begin{tabularx}{\textwidth}{|>{\raggedright\arraybackslash}p{3.1cm}|X|}
\hline
\textbf{Field} & \textbf{Description} \\
\hline
Test ID & IT-02 \\
\hline
Objective & \textit{read\_telemetry()} returns fresh distances at the firmware’s print rate when the loop runs. \\
\hline
Input & Powered robot with sensors unobstructed and obstructed. \\
\hline
Procedure & Run \textit{python3 uart\_bridge.py} or full \textit{main.py}; log lines. \\
\hline
Expected & Non-\textit{None} dicts when the line format matches; values change when obstacles move. \\
\hline
Observed & Operator records sample log excerpt for the report appendix if required. \\
\hline
\end{tabularx}
\end{table}

%-------------------------------------------------------------------
\section{System Testing}
\label{sec:system_testing}

System tests evaluate end-to-end behaviour: sensing, inference (when online), validation, and actuation. \textbf{Tables~\ref{tab:st_auto}},~\textbf{\ref{tab:st_ai}}, and~\textbf{\ref{tab:st_safe}} document the system test procedures.

\subsection[AUTO mode avoidance]{AUTO\_MODE obstacle avoidance}
\label{subsec:st_auto}

\begin{table}[H]
\centering
\caption[System test: standalone firmware behaviour]{\textbf{System test: standalone firmware behaviour}}
\label{tab:st_auto}
\begin{tabularx}{\textwidth}{|>{\raggedright\arraybackslash}p{3.1cm}|X|}
\hline
\textbf{Field} & \textbf{Description} \\
\hline
Test ID & ST-01 \\
\hline
Objective & With Pi silent, robot executes rule-based avoidance using thresholds (e.g.\ \textit{OBSTACLE\_DIST}, \textit{WALL\_DIST}). \\
\hline
Input & Cluttered tabletop course; Pi not sending JSON. \\
\hline
Procedure & Power on; allow AUTO\_MODE; observe detours and stops. \\
\hline
Expected & No continuous forward into reported obstacle; recovery without watchdog fault. \\
\hline
Observed & Qualitative pass/fail; note any tuning changes to thresholds. \\
\hline
\end{tabularx}
\end{table}

\subsection[AI mode and fallback]{AI\_MODE planning and UART timeout fallback}
\label{subsec:st_ai}

\begin{table}[H]
\centering
\caption[System test: intelligence layer and mode fallback]{\textbf{System test: intelligence layer and mode fallback}}
\label{tab:st_ai}
\begin{tabularx}{\textwidth}{|>{\raggedright\arraybackslash}p{3.1cm}|X|}
\hline
\textbf{Field} & \textbf{Description} \\
\hline
Test ID & ST-02 \\
\hline
Objective & Valid JSON from Pi drives motion; silence beyond \textit{AI\_TIMEOUT\_MS} restores AUTO\_MODE~\cite{brooks1986robust}. \\
\hline
Input & Running \textit{main.py} with valid API key; then stop Pi process or disconnect TX. \\
\hline
Procedure & Observe OLED/mode indicators and behaviour before and after timeout. \\
\hline
Expected & Initial AI\_MODE motion consistent with validator output; after timeout, autonomous avoidance resumes. \\
\hline
Observed & Record timeout duration and qualitative result. \\
\hline
\end{tabularx}
\end{table}

\subsection[Safe stop on faults]{Safe stop on inference or serial faults}
\label{subsec:st_safe}

\begin{table}[H]
\centering
\caption[System test: degraded operation]{\textbf{System test: degraded operation}}
\label{tab:st_safe}
\begin{tabularx}{\textwidth}{|>{\raggedright\arraybackslash}p{3.1cm}|X|}
\hline
\textbf{Field} & \textbf{Description} \\
\hline
Test ID & ST-03 \\
\hline
Objective & Orchestrator commands \textit{stop} when telemetry is missing, JSON is unparseable, or the LLM call raises~\cite{vemprala2024chatgpt}. \\
\hline
Input & Block sensors (optional), revoke key temporarily, or disconnect UART mid-run. \\
\hline
Procedure & Trigger each fault class once; capture logs. \\
\hline
Expected & No sustained forward motion under fault; log lines show stop or skip with warning. \\
\hline
Observed & Document per scenario. \\
\hline
\end{tabularx}
\end{table}

%-------------------------------------------------------------------
\section{Requirements Traceability}
\label{sec:traceability}

\textbf{Table~\ref{tab:trace}} maps selected functional requirements from Chapter~\ref{ch:srs} to the tests above. Full coverage of FR-01--FR-15 is completed during project examination by extending the same table format.

\begin{table}[H]
\centering
\caption[Partial requirements-to-test mapping]{\textbf{Partial requirements-to-test mapping}}
\label{tab:trace}
\begin{tabularx}{\textwidth}{|>{\raggedright\arraybackslash}p{2.4cm}|X|>{\raggedright\arraybackslash}p{2.8cm}|}
\hline
\textbf{Requirement} & \textbf{Intent} & \textbf{Tests} \\
\hline
FR-06, FR-07 & Validation and safe stop & UT-01 \\
\hline
FR-02, FR-03 & Telemetry and parsing & UT-02, IT-02 \\
\hline
FR-08, FR-09 & JSON command on UART & IT-01 \\
\hline
FR-12, FR-13 & Modes and fallback & ST-01, ST-02 \\
\hline
FR-15 & Fault logging and stop & ST-03 \\
\hline
\end{tabularx}
\end{table}

%-------------------------------------------------------------------
% \section{Summary}
% \label{sec:ch7summary}

% This chapter defined a three-level test strategy for KANDA: unit checks on the validator and telemetry parser, integration checks on the UART link, and system checks for AUTO\_MODE, AI\_MODE, timeout fallback, and safe degradation. Tabulated procedures support repeatability during demonstration and viva. Limitations include dependence on manual observation for motion quality and on network availability for cloud-backed ST-02; future work may add scripted hardware-in-the-loop logging and timestamped test reports.
